\documentclass[12pt,a4paper]{report}

% #region Packages
\renewcommand{\rmdefault}{phv}	% Utilise la police Arial pour le texte
\usepackage[T1]{fontenc}		% Encodage de la police pour une meilleure gestion des caractères spéciaux
\usepackage[utf8]{inputenc}		% Encodage UTF-8 pour permettre l'utilisation de caractères non-ASCII
\usepackage{lmodern}			% Utilise une version améliorée des polices LaTeX standard
\usepackage{graphicx}			% Permet d'insérer des images dans le document
\usepackage{multicol}			% Gère le texte en colonnes multiples
\usepackage{lipsum}				% Génère du texte factice (Lorem Ipsum)
\usepackage{float}				% Améliore le placement des objets flottants (comme les figures et tableaux)
\usepackage{vwcol}				% Permet de créer des colonnes avec des largeurs variables
\usepackage{geometry}			% Permet de définir les marges et la mise en page
\usepackage{caption}			% Personnalise les légendes des figures et tableaux
\usepackage{framed}				% Encadre des portions de texte
\usepackage{pdflscape}			% Permet d'inclure des pages en mode paysage dans un PDF
\usepackage{eurosym}			% Ajoute le symbole de l'euro (€)
\usepackage{xcolor}				% Gestion avancée des couleurs dans le document
\usepackage{amsmath,amssymb}	% Fournit des symboles et des outils avancés pour les mathématiques
\usepackage{hyperref}			% Ajoute des liens hypertextes cliquables dans le PDF
\usepackage{setspace}			% Gère l'espacement entre les lignes

% #endregion

% #region Déclaration des couleurs
\definecolor{tabblue}{RGB}{31,119,180}
\definecolor{taborange}{RGB}{255,127,14}
\definecolor{tabgreen}{RGB}{44,160,44}
\definecolor{tabred}{RGB}{214,39,40}
\definecolor{tabpurple}{RGB}{148,103,189}
\definecolor{tabbrown}{RGB}{140,86,75}
\definecolor{tabpink}{RGB}{227,119,194}
\definecolor{tabgray}{RGB}{127,127,127}
\definecolor{tabolive}{RGB}{188,189,34}
\definecolor{tabcyan}{RGB}{23,190,207}
\definecolor{rosePW}{RGB}{219,164,207}
\definecolor{orangeNW}{RGB}{245,121,0}
% #endregion

\geometry{a4paper, left=25mm, right=25mm, top=25mm, bottom=25mm}
\DeclareCaptionType{annexe}


% Espacement entre les lignes
\onehalfspacing



\usepackage{hyperref}
\hypersetup{
    colorlinks=true,
    linkcolor=black,
    filecolor=black,
    urlcolor=black,
    % bookmarks=true,
    bookmarksopen=true,
    bookmarksopenlevel=2,
    pdftitle={Manuscrit de thèse de Maxime Guillot},
    pdfauthor={Maxime Guillot},
}

% Set up section numbering
\setcounter{secnumdepth}{4}


\usepackage{titlesec}
\titleformat{\paragraph}
{\normalfont\normalsize\bfseries}{\theparagraph}{1em}{}
\titlespacing*{\paragraph}
{0pt}{3.25ex plus 1ex minus .2ex}{1.5ex plus .2ex}


% #region Déclaration des "Mets une source"
\usepackage{etoolbox}
\BeforeBeginEnvironment{metsUneSource}{\colorlet{oldcolor}{.}}
\AfterEndEnvironment{metsUneSource}{\color{oldcolor}}

\newenvironment{metsUneSource}
{\begin{framed}\color{red}}
{\end{framed}}
% #endregion

% % Création de l'environnement 'spaced' avec interligne doublé
% \newenvironment{spaced}
%     {\begin{spacing}{2}} % Interligne doublé
%     {\end{spacing}}

\newcommand{\HRule}{\rule{\linewidth}{0.5mm}} % Une ligne horizontale de la largeur de la page








\begin{document}
% ----------------------------- %
%          Document principal    %
% ----------------------------- %

% Inclusion de la première de couverture
\begin{titlepage}
    \begin{center}
        % Logo de l'université (facultatif)
        \includegraphics[width=0.2\textwidth]{figures/universite_de_bordeaux.pdf} \\
        \vspace{1cm}

        % Nom de l'université
        \textsc{\Large Université de [Nom de ton université]}\\[0.5cm]

        % Faculté ou département
        \textsc{\Large [Nom de la faculté ou du département]}\\[1.5cm]
\begin{spacing}{2}
     % Titre de la thèse
     \HRule \\[0.4cm]
     { \huge \bfseries Référence de tension en technologie CMOS désensibilisée au débit de dose en environnement radiatif sévère : un verrou technologique pour le démantèlement des centrales nucléaires}\\[0.4cm]
     \HRule \\[1.5cm]
\end{spacing}
       

        % Sous-titre (si applicable)
        { \large Voltage reference in CMOS technology desensitized to the dose rate in a harsh radiation environment for the dismantling of nuclear power plants}\\[1.5cm]

        % Auteur
        \textsc{\Large Présenté par :}\\[0.4cm]
        \textsc{\large Maxime [Ton nom de famille]}\\[2cm]

        % Diplôme
        \textsc{\Large En vue de l'obtention du diplôme de [Nom du diplôme]}\\[1cm]

        \begin{flushleft}
            % Jury
            \textsc{Jury :}\\[0.4cm]
            \begin{tabular}{ll}
                Prénom NOM & Fonction \\
                Prénom NOM & Fonction \\
                Prénom NOM & Fonction \\
                Prénom NOM & Fonction \\
                Prénom NOM & Fonction \\
            \end{tabular}
        \end{flushleft}
 
    \end{center}
\end{titlepage}




% ----------------------------- %
%          Préface               %
% ----------------------------- %
\chapter*{Préface}
\addcontentsline{toc}{chapter}{Préface}  % Ajout au sommaire
% Écrire la préface ici

% ----------------------------- %
%        Table des matières      %
% ----------------------------- %
\tableofcontents
\newpage

% ----------------------------- %
%      Chapitre d'introduction   %
% ----------------------------- %
\chapter{Introduction}
% Écrire l'introduction ici
 

% ----------------------------- %
%        Chapitre 1              %
% ----------------------------- %
\chapter{Titre du Chapitre 1}
% Inclure un fichier séparé pour chaque chapitre
% \input{chapitre1}

metsUneSource


\lipsum[1-2]

\begin{metsUneSource}
\lipsum[3]
\end{metsUneSource}

\textit{dcnduvbdvbdbvdiubvk}

% ----------------------------- %
%        Chapitre 2              %
% ----------------------------- %
\chapter{Titre du Chapitre 2}
% Inclure un fichier séparé pour chaque chapitre
% \input{chapitre2}

% ----------------------------- %
%       Conclusion               %
% ----------------------------- %
\chapter{Conclusion}
% Écrire la conclusion ici

% ----------------------------- %
%      Bibliographie             %
% ----------------------------- %
% \addcontentsline{toc}{chapter}{Bibliographie}
% \bibliographystyle{plain}  % Style de la bibliographie
% \bibliography{bibliographie}  % Fichier .bib

% \bibliographystyle{IEEEtran}
% \bibliography{MaBibliographie}


\begin{titlepage}
    \vspace*{\fill}
    \begin{center}
        {\large \textsc{Résumé de la thèse}} \\[1cm]
        % Résumé de la thèse (environ 150-300 mots)
        \noindent Cette thèse traite de [sujet principal] et explore [principaux thèmes ou contributions]. L'objectif principal est de [résumé des objectifs]. Les résultats obtenus montrent que [résumé des résultats].\\[2cm]

        % % Mots-clés
        % \textsc{\large Mots-clés :}\\
        % [Mots-clés 1, Mots-clés 2, Mots-clés 3, etc.]

        % Mentions légales (facultatif)
        \textsc{\small [Informations légales ou institutionnelles si nécessaire]}\\
    \end{center}
    \vspace*{\fill}
\end{titlepage}


\end{document}
